\documentclass{sprawozdanie-agh}

\usepackage[utf8]{inputenc}
\usepackage{matlab-prettifier}
\usepackage{siunitx}
\usepackage{booktabs,caption,ragged2e}
\usepackage[shortlabels]{enumitem}


\makeatletter

% ========= Początek Dokumentu ========= 

\begin{document}
% ========= Strona tytułowa ========= 
\przedmiot{Teoria Sterowania II: Ćwiczenia Laboratoryjne}
\tytul{Wytyczne dla Ćwiczenia 3}
\podtytul{Metoda Lapunowa}
\kierunek{Automatyke i~Robotyka}
\autor{Dariusz Cieślar, Maciej Różewicz}
\data{Kraków, 3 marca 2026}

\stronatytulowa{}

% ========= Treść Właściwa =========
\section{Cel ćwiczenia}
Celem ćwiczenia jest zapoznanie się z~metodami badania stabilności układów
nieliniowych.
W czasie zajęć przedstawiona i~przećwiczona zostanie pośrednia metodą Lapunowa
badania stabilności punktów równowagi układów nieliniowych.

Jest to stosunkowo prosta w~zastosowaniu metoda, dająca zadowalające wyniki
w~wielu praktycznych przypadkach, jednak mająca pewne ograniczenia.


%%%%%%%%%%%%%%%%%%%%%%%%%%%%%%%%%%%%%%%%%%%%%%%%%%%%%%%%%%%% 
\section{Przebieg ćwiczenia}
    %=======================================================
    \subsection{Pośrednia metoda Lapunowa}
    Dla każdego z~podanych systemów opisanych w podpunktach od \ref{sec:posrednia_1} do \ref{sec:posrednia_5} należy:
    \begin{enumerate}[label=\textbf{Krok \arabic*}]
        \item Wyznaczyć punkty równowagi układu.
        \item Wyznaczyć macierz Jacobiego $A = \frac{\partial f}{\partial x}\bigg|_{x=x_e}$ układu w~każdym z wyznaczonych punktów równowagi.
        %\item Zamodelować układ równań w~Simulinku (lub użyć skryptu, np. funkcji \textit{ode45}, do jego rozwiązania).
        \item Wyznaczyć portrety fazowe układu nieliniowego.
        \item Wyznaczyć portrety fazowe linearyzacji układu w~każdym z~punktów równowagi. 
        \item Sprawdzić, czy można zastosować pośrednią metodę Lapunowa do badania stabilności punktów równowagi.
        \item Zbadać stabilność wszystkich punktów równowagi, dla których jest to możliwe.
    \end{enumerate}
        %---------------------------------------------------
        \subsubsection{System 1} \label{sec:posrednia_1}
        Rozważamy układ mechaniczny masa-sprężyna ze sprężyną o~nieliniowej charakterystyce.
        Układ ten opisany jest równaniem (\ref{eq:posrednia_1}).
        \begin{equation}\label{eq:posrednia_1}
            \ddot{x} + b\dot{x} + cx + dx^{3} = 0, \, b,c > 0,\, |c|<|d|.
        \end{equation}
        Gdzie:
        \begin{itemize}
            \item $x$ - odchylenie od punktu równowagi,
            \item $b$ - współczynnik tarcia,
            \item $c$ i~$d$ - parametry sprężyny.
        \end{itemize}
        Proszę rozważyć przypadek dla $d>0$ i~$d<0$.

        %---------------------------------------------------
        \subsubsection{System 2} \label{sec:posrednia_2}
        Rozważamy wahadło z~tłumieniem, opisywane równaniem:
        \begin{equation}\label{eq:posrednia_2}
            \ddot{x} + \frac{c}{lm}\dot{x} + \frac{g}{l}\sin{x} = 0
        \end{equation}
        gdzie:
        \begin{itemize}
            \item $x$ - kąt wychylenia wahadła z~punktu równowagi trwałej,
            \item $g$ - przyspieszenie ziemskie,
            \item $m$ - masa wahadła,
            \item $l$ - długość wahadła.
        \end{itemize}

        %---------------------------------------------------
        \subsubsection{System 3} \label{sec:posrednia_3}
        Rozważamy układ równań Van der Pola:
        \begin{equation}
            \label{eq:posrednia_3}
            \left\{
            \begin{array}{l}
                 \dot{x}_{1} = x_{2} - x_{1}^{3} - ax_{1}\\
                 \dot{x}_{2} = -x_{1}
            \end{array}
            \right\}.
        \end{equation}
        Proszę zbadać zachowanie układu dla różnych wartości $a$.

        %---------------------------------------------------
        \subsubsection{System 4} \label{sec:posrednia_4}
        Rozważamy układ równań wahadła bez tłumienia:
        \begin{equation}
            \label{eq:posrednia_4}
            \left\{
            \begin{array}{l}
                 \dot{x}_{1} = x_{2}\\
                 \dot{x}_{2} = -\sin{x_1}
            \end{array}
            \right\}.
        \end{equation}

        %---------------------------------------------------
        \subsubsection{System 5} \label{sec:posrednia_5}
        Rozważamy układ równań modelujący konkurencję dwóch gatunków:
        \begin{equation}
            \label{eq:posrednia_5}
            \left\{
            \begin{array}{l}
                 \dot{x}_{1} = x_2\\
                 \dot{x}_{2} = x_1 - x_1^2 - x_2
            \end{array}
            \right\}.
        \end{equation}


%%%%%%%%%%%%%%%%%%%%%%%%%%%%%%%%%%%%%%%%%%%%%%%%%%%%%%%%%%%% 
\section{Przebieg ćwiczenia}
    Przygotuj sprawozdanie z ćwiczenia opisujące cel, przebieg i obserwacje.
    
    W części opisującej przebieg ćwiczenia proszę skupić się na przedstawieniu
    wyników dla każdego z~badanych układów, wraz z~odpowiednimi portretami
    fazowymi i~analizą stabilności punktów równowagi.

    W sprawozdaniu muszą znaleźć się:
    \begin{itemize}
        \item wyznaczone punkty równowagi,
        \item wyznaczone macierze Jacobiego,
        \item portrety fazowe układu nieliniowego,
        \item portrety fazowe linearyzacji układu,
        \item analiza możliwości zastosowania pośredniej metody Lapunowa,
        \item analiza stabilności punktów równowagi.
    \end{itemize}
\end{document}
